\documentclass[main.tex]{subfiles}

\begin{document}

\section{Future discussions}
\begin{frame}
  \frametitle{Specialized kernels}
  \small
  A general SGEMM kernel must handle many shapes and layouts. Real models let
  us trade some of that flexibility for more performance.

  \begin{columns}[T,onlytextwidth]
    \begin{column}{0.47\textwidth}
      \begin{block}{What can be specialized?}
        \begin{itemize}
          \item Fixed tensor shapes and memory layouts
          \item Reduced precision and Tensor Core tile sizes
          \item Structured sparsity or quantized weights
          \item Fused bias, activation, and normalization
        \end{itemize}
      \end{block}
    \end{column}
    \begin{column}{0.49\textwidth}
      \begin{block}{Why does it help?}
        \begin{itemize}
          \item Less indexing and control-flow overhead
          \item Fewer kernel launches and intermediate writes
          \item Better use of registers and specialized hardware
          \item Autotuning can choose a kernel for each layer
        \end{itemize}
      \end{block}
    \end{column}
  \end{columns}

  \begin{alertblock}{Example: fuse the epilogue}
    Compute \(Y=\operatorname{ReLU}(XW+b)\) in one kernel instead of writing
    \(XW\) to global memory and reading it back for bias and activation.
  \end{alertblock}
\end{frame}

\begin{frame}
  \frametitle{Continuing GPU evolution}
  \only<1>{
    \begin{columns}[T,onlytextwidth]
      \begin{column}{0.42\textwidth}
        \textbf{2017: Volta Tensor Cores}
        \begin{itemize}
          \item Tensor Cores perform mixed-precision matrix multiply--accumulate:
            \(D=A B+C\).
          \item The original Volta unit multiplies FP16 inputs and accumulates in
            FP16 or FP32.
          \item Matrix multiplication becomes an explicit hardware primitive,
            alongside general-purpose CUDA cores.
        \end{itemize}
      \end{column}
      \begin{column}{0.56\textwidth}
        \centering
        \begin{tikzpicture}[font=\scriptsize,>=Latex]
  \node[draw,fill=blue!20,minimum size=1.1cm] (a) at (0,1) {\(A_{4\times4}\)};
  \node at (.85,1) {\(\times\)};
  \node[draw,fill=green!20,minimum size=1.1cm] (b) at (1.7,1) {\(B_{4\times4}\)};
  \node at (2.55,1) {\(+\)};
  \node[draw,fill=orange!20,minimum size=1.1cm] (c) at (3.4,1) {\(C_{4\times4}\)};
  \node[draw,rounded corners,fill=purple!18,minimum width=4.5cm,minimum height=.75cm] (tc) at (1.7,-.05) {Tensor Core: matrix multiply--accumulate};
  \draw[-Latex] (a)--(tc);
  \draw[-Latex] (b)--(tc);
  \draw[-Latex] (c)--(tc);
  \node[draw,fill=red!15,minimum width=1.2cm,minimum height=.65cm] (d) at (1.7,-1.15) {\(D_{4\times4}\)};
  \draw[-Latex] (tc)--(d);
\end{tikzpicture}


        {\tiny Operation shape shown for first-generation Volta Tensor Cores. Source: NVIDIA, \emph{Volta Architecture Whitepaper}, 2017.}
      \end{column}
    \end{columns}
  }

  \only<2>{
    \vspace{-2pt}
    \begin{columns}[T,onlytextwidth]
      \begin{column}{0.44\textwidth}
        \textbf{2020 onward: specialize around AI dataflow}
        \begin{itemize}
          \item Ampere (2020) adds TF32 and structured sparsity support.
          \item Hopper (2022) adds FP8 Tensor Cores and a Transformer Engine that
            dynamically chooses FP8 or FP16 per layer.
          \item Blackwell (2024) adds more low-precision formats and a second-generation
            Transformer Engine. This is more than incremental scaling.
        \end{itemize}
      \end{column}
      \begin{column}{0.54\textwidth}
        \centering
        \begin{tikzpicture}[font=\scriptsize,node distance=4mm,>=Latex]
  \node[draw,rounded corners,fill=blue!15,minimum width=3.8cm] (layer) {Transformer layer};
  \node[draw,rounded corners,fill=orange!25,below=of layer,minimum width=3.8cm,align=center] (te) {Transformer Engine\\select scale + FP8/FP16};
  \node[draw,rounded corners,fill=green!20,below=of te,minimum width=3.8cm,align=center] (tc) {Tensor Cores\\matrix multiply};
  \node[draw,rounded corners,fill=purple!15,below=of tc,minimum width=3.8cm] (out) {next layer};
  \draw[-Latex] (layer)--(te);
  \draw[-Latex] (te)--(tc);
  \draw[-Latex] (tc)--(out);
\end{tikzpicture}


        {\tiny Source: NVIDIA, \emph{Hopper Architecture In-Depth} (2022).}
      \end{column}
    \end{columns}
  }

  \only<3>{
    \textbf{2026: Vera Rubin---optimization moves to the rack}

    \vspace{1pt}
    \begin{columns}[T,onlytextwidth]
      \begin{column}{0.59\textwidth}
        \begin{itemize}
          \item The target is no longer an isolated GPU, but a complete
            rack-scale computing platform.
          \item Vera CPUs, Rubin GPUs, memory, NVLink, and networking are
            designed together to run one distributed AI workload.
          \item A high-bandwidth scale-up fabric lets many accelerators share
            model parameters and intermediate results.
          \item Performance now depends on communication, power, cooling, and
            software orchestration as well as matrix throughput.
        \end{itemize}
      \end{column}
      \begin{column}{0.37\textwidth}
        \begin{block}{Unit of optimization}
          \centering
          Kernel\\[3pt]
          $\downarrow$\\[3pt]
          GPU\\[3pt]
          $\downarrow$\\[3pt]
          Server\\[3pt]
          $\downarrow$\\[3pt]
          \textbf{Rack-scale system}
        \end{block}
      \end{column}
    \end{columns}

    \vfill
    {\tiny Source: NVIDIA, \emph{Vera Rubin Architecture}, announced roadmap.}
  }

\end{frame}

\begin{frame}
  \frametitle{Alternative accelerator architectures}
  \small

  \only<1>{
    \textbf{Different vendor, or different architecture?}
    \vspace{3pt}
    \begin{columns}[T,onlytextwidth]
      \begin{column}{0.48\textwidth}
        \begin{block}{AMD Instinct: another GPU ecosystem}
          AMD's datacenter GPUs use many programmable compute units, matrix
          instructions, and high-bandwidth memory. ROCm provides the software
          stack used to program them.

          \medskip
          \textbf{Strength:} retains GPU flexibility while competing with CUDA.

          \textbf{Tradeoff:} performance depends heavily on software and library
          support, not only the silicon.
        \end{block}
      \end{column}
      \begin{column}{0.48\textwidth}
        \begin{block}{Google TPU: systolic arrays}
          Data moves rhythmically through a two-dimensional grid of
          multiply--accumulate units. Values are reused as they pass between
          neighboring units instead of repeatedly returning to memory.

          \medskip
          \textbf{Strength:} efficient, predictable dense tensor operations.

          \textbf{Tradeoff:} less flexible for irregular workloads.
        \end{block}
      \end{column}
    \end{columns}

    \vfill
    {\tiny Sources: AMD, \emph{Instinct Accelerators}; Jouppi et al. (2017).}
  }

  \only<2>{
    \textbf{Other designs change the scale or location of computation.}
    \vspace{3pt}
    \begin{columns}[T,onlytextwidth]
      \begin{column}{0.48\textwidth}
        \begin{block}{Cerebras: wafer-scale integration}
          Many compute cores, distributed SRAM, and an on-wafer fabric are
          placed on one large processor. More communication stays on the wafer
          instead of crossing between accelerator packages.

          \medskip
          \textbf{Strength:} memory capacity and bandwidth close to compute.

          \textbf{Tradeoff:} difficult manufacturing, cooling, and mapping.
        \end{block}
      \end{column}
      \begin{column}{0.48\textwidth}
        \begin{block}{NPUs: AI inside the device}
          NPUs sit beside CPUs and GPUs in phones, laptops, and embedded
          systems. They accelerate supported quantized tensor operations under
          strict power and thermal limits.

          \medskip
          \textbf{Strength:} private, low-latency inference without the cloud.

          \textbf{Tradeoff:} limited memory and operator support.
        \end{block}
      \end{column}
    \end{columns}

    \vfill
    {\tiny Sources: Cerebras, \emph{Wafer-Scale Engine}; Sze et al. (2017).}
  }

  \only<3>{
    \begin{columns}[T,onlytextwidth]
      \begin{column}{0.52\textwidth}
        \begin{block}{Model ASIC: Etched Sohu}
          The hardware and dataflow target transformer inference rather than
          arbitrary kernels. Removing general-purpose machinery can devote more
          silicon and energy to the intended model family.

          \medskip
          \textbf{Strength:} potentially much higher throughput per watt.

          \textbf{Tradeoff:} model changes can make fixed hardware obsolete.
        \end{block}
      \end{column}
      \begin{column}{0.44\textwidth}
        \begin{block}{The design spectrum}
          \centering
          GPU: programmable\\[3pt]
          $\downarrow$\\[-1pt]
          TPU/NPU: operator-focused\\[3pt]
          $\downarrow$\\[-1pt]
          ASIC: model-focused
        \end{block}
        \centering
        \scriptsize More specialization $\Rightarrow$ less flexibility
      \end{column}
    \end{columns}

    \begin{alertblock}{There is no universally best architecture}
      Specialization can improve performance per watt, but reduces flexibility.
      The right balance depends on workload stability, scale, power, and the
      cost of moving data.
    \end{alertblock}

    \vfill
    {\tiny Source: Etched (2024), ``Introducing Sohu.''}
  }
\end{frame}

\end{document}
